% Isolated Version 14 author draft; assembled independent review is required.
% This file does not modify the original polynomial or Version 13.
\section{A quartic presentation with 299 variables}\label{sec:v14}

We retain the arithmetic inputs, qualifications and source editions of
Sections~\ref{subsec:v12-inputs} and~\ref{sec:v13}. The polynomial of
the main construction is unchanged. The following are smaller
presentations of the same rational input projection.
\begin{theorem}\label{thm:v14-quartic}
There is a specified integer polynomial $P_{299}$ of total degree four
in 299 variables such that
\[
 X\in\mathbb Z\quad\Longleftrightarrow\quad
 \exists\mathbf v\in\mathbb Q^{298}:P_{299}(X,\mathbf v)=0.
\]
It has 6,630 nonzero monomials, content one, and largest absolute
coefficient $186624\cdot3^{2^{26}}$. Its coefficient artifact writes
$C=3^{2^{25}}$ as a fixed integer, not a quantified variable.
There is also a specified integer quartic $P_{314}$ with
\[
 X\in\mathbb Z\quad\Longleftrightarrow\quad
 \exists\mathbf w\in\mathbb Q^{313}:P_{314}(X,\mathbf w)=0.
\]
It has 6,672 nonzero monomials and content one. Every coefficient has
absolute value strictly less than $2^{6496}$; its largest coefficient
has exactly 6,496 bits. The compact ordinary decimal coefficient table
occupies 264,430 bytes in its specified ASCII format. The separate table
retaining source coordinate identifiers occupies 222,890 bytes.
\end{theorem}
The ordinary coefficient bound here is different from the bound for
$P_{411}$. We prove the new bound for the new artifact. Some changes
below are inverse changes of private witnesses; others change an
arithmetic certificate or point domain. For the latter we prove
soundness and completeness separately.

\subsection{Conventions and retained interfaces}\label{subsec:v14-inputs}
Put $K=\mathbb Q(\rho)$, $\rho^2=2$, $d=3-\rho$, $D=d^2=11-6\rho$,
and $E:y^2=x^3-Dx$. Retain
$E(K)=\mathbb ZR_0+E[2](K)$ and $2E(K)=\mathbb ZP$, where
\[
 R_0=(1+\rho,2\rho),\quad P=(c_*,a_*b_*),\quad
 a_*=-1+7\rho/4,\quad b_*=39/8-9\rho/2,\quad c_*=a_*^2.
\]
Thus $b_*^2=c_*^2-D$. Write $f(Q)=1/x(Q)$ and
$z(Q)=-x(Q)/y(Q)$, with both functions extended regularly by zero
at $O$. The roots are $Q_i=[n_i]P$ with labels
$(n_1,n_2,n_3,n_4)=(n,a,t,b)$. Their nonzero odd integral labels
will be restored before height or logarithm arguments. The scalar
point is $Q_2-P=[s]P$, where $s=a-1$.

The retained inputs include the squareclass calculation, formal
expansions and logarithm estimates, simultaneous contact-prime
construction with its real-density condition, preparation-zero pole
criterion, root-height implication, Lucas argument and scalar-depth-24
argument. In particular the three positive rational square scales are
\[
 A_0=1,\quad A_1=C_0^2,\quad A_2=(1+C_1^2)^2,\qquad
 \prod_i x(Q_i)=C_0+C_1\rho .
\]
The same root-height conclusion is
\[
 B^2-6<7\mathcal H|N|^{19/k},\qquad
 B=\max_i|n_i|,\quad 1\leq\mathcal H<2^{2093966},
\]
when the $k$th powers of the three required contact radicals divide
the nonzero integer $N$. The independently restored scalar anchor
gives $B>(40\mathcal H)^{10}$, excluding all three radicals equal
to one without using H or auxiliary logarithms.
The retained preparation-zero equations exclude negative odd
valuations of a pole input at required large primes. Their
completeness interface uses the same small-prime squareclasses and
depths, square reduced denominator, and real interval $(0,17)$.
We verify the changed consumers of these inputs below.

\subsection{The fixed-difference chart for H}\label{subsec:v14-H}
First replace the old $x(H)$ pair by the equations
\begin{equation}\label{eq:v14-H-raw}
 c_*(f+F)+Dv^2+2b_*v=1,\qquad v=ft,\qquad F=vt.
\end{equation}
The common primary normalizer is the fixed affine expression
\begin{equation}\label{eq:v14-normalizer}
 \mathcal S=(b_*+c_*t+Dv)/a_*.
\end{equation}
Its substitution into every primary, logarithm and scalar comparison
row preserves quadratic degree. If $f=0$, the last two equations
force $v=F=0$, contradicting the first. Set $x=1/f$. They give
$v^2=fF$ and $a_*f\mathcal S-b_*f-(c_*+Df)v=0$.
Here is an unconditional certificate for the recovered curve.
Put $d_1=c_*+Df$, $B_1=b_*f+d_1v$,
$A_1=a_*f\mathcal S-B_1$, and
$K_1=c_*(f+F)+Dv^2+2b_*v-1$. Then
\[
 c_*f(f\mathcal S^2+Df^2-1)
 =fd_1K_1+c_*d_1(v^2-fF)
 +(b_*^2-c_*^2+D)f^2+A_1(a_*f\mathcal S+B_1).
\]
Thus $\mathcal S^2=x-Df$ and the old fixed-difference equations
hold. This direction divides only by $c_*f$. Conversely, the old
full chart gives $t=v/f$ and the same affine expression for
$\mathcal S$; its previously established exceptional-point exclusion
supplies the forward implication. All outside values are preserved.

Next use the already needed outputs $U=f(H)^2$, $V=f(H-P)^2$.
Store $W,v,t,U,V$, put $Q=(1-DW-2b_*v)/c_*$, and impose
\begin{equation}\label{eq:v14-H-squared}
 W=v^2,\qquad UV=W^2,\qquad Q^2=U+V+2W,\qquad t(U+W)=Qv.
\end{equation}
On $UQ\ne0$, set $f=(U+W)/Q$, $F=(V+W)/Q$. If
$E_1=UV-W^2$ and $E_2=Q^2-U-V-2W$, then
\begin{align*}
 (U+W)^2-UQ^2&=-UE_2-E_1,\\
 (V+W)^2-VQ^2&=-VE_2-E_1,\\
 (U+W)(V+W)-WQ^2&=-WE_2+E_1.
\end{align*}
These give $f^2=U$, $F^2=V$, $fF=W$, $f+F=Q$, $ft=v$,
and $F=vt$, restoring~\eqref{eq:v14-H-raw}.

The extra branches are classified before this reconstruction.
For $U=0$ one has $W=v=0$, $Q=1/c_*$, $V=1/c_*^2$.
For $Q=0$ one has $U=V=-v^2$ and
\[
 Dv^2+2b_*v-1=0,\qquad
 v=\frac{3+2\rho}{4}\quad\hbox{or}\quad
 v=\frac{-36+16\rho}{49}.
\]
These are units at every place above $p>76176$:
$c_*+b_*=12-8\rho$ has norm16 and
$c_*-b_*=(9+4\rho)/4$ has norm$49/16$.
Section~\ref{subsec:v14-groups} excludes these branches by their
unchanged $V_0$ unit leaf and the pre-H height/anchor contradiction.
No division by $UQ$ is used to prove its own legitimacy.

\subsection{Root parity and the shifted anchor}\label{subsec:v14-parity}
The root chart $c_i=a_i^2$, $b_i^2=c_i^2-D$ represents
$(c_i,a_ib_i)$. Its abscissa is nonzero; for the final presentation
this follows from the retained product-height inverse as proved in
Section~\ref{subsec:v14-reciprocal}.
The squareclass map has independent values
\[
 \alpha(R_0)=[1+\rho],\qquad \alpha(T_0)=[-1],\qquad
 \alpha(T_d)=[d].
\]
Indeed norms in a putative square relation force the exponents
of $1+\rho$ and $d$ even, and the real embeddings exclude the
remaining class $[-1]$. Thus the square abscissa already puts
each root in $2E(K)=\mathbb ZP$, with nonzero integral label,
before an odd-chart condition.

Keep the scalar reciprocal equations
\[
 (c_*c_2+D)g=a_*b_2-b_*a_2,\qquad
 g=g_0+g_1\rho,\qquad g_0\tau=1.
\]
The denominator is nonzero at both real embeddings. The subtraction
identity gives $f(Q_2-P)=g^2$. At equal abscissae, $Q_2=P$
contradicts the guard and $Q_2=-P$ gives the exact duplication
value for $-2P$. Hence $s=n_2-1\ne0$ is an integer before height
or logarithms. Write $F_s=2g_0g_1$.

Let $X_n$ be the $\rho$ coefficient of $x([n]P)$. For $n\ne0$,
\begin{equation}\label{eq:v14-X}
 v_3(X_n)=
 \begin{cases}0&n\text{ odd},\\-2-2v_3(n)&n\text{ even},\end{cases}
 \qquad v_2(X_n)=-1\quad(n\text{ odd}).
\end{equation}
At3 the reduction of $P$ is $(\rho,0)$ over $\mathbb F_9$,
which proves the odd case. For $n=2\cdot3^ku$, $u$ a3-unit,
the formal expansion gives
$z([n]P)/3^{k+1}\equiv u(2+\rho)\pmod3$.
Since $x(z)=z^{-2}+O(z^2)$ and
$(2+\rho)^{-2}=3/2-\rho\equiv2\rho\pmod3$, the scaled
abscissa has nonzero imaginary coefficient modulo3. This proves
the even case and $X_n\ne0$.
At2, $z(P)=(232+86\rho)/357$ has component valuations3 and1.
The integral odd multiplication series preserves these values:
its error has field valuation at least$9/2$ when $v(2)=1$.
For $z=A+B\rho$ with those valuations the imaginary coefficient
of $z^{-2}$ is $-2AB/(A^2-2B^2)^2$, of valuation$-1$.
The Laurent correction has larger valuation. This uses integral
formal coefficients, not good reduction at2.

The retained scalar-coefficient lemma gives $v_3(F_n)=0$ for odd
$n$, $2+2v_3(n)$ for even nonzero $n$, and dyadic integrality.
It is applied with both parities. Put
\begin{equation}\label{eq:v14-q}
 q=12(X_A^2+X_T^2)+108X_Z^2,\qquad T=F_s+Cq,
\end{equation}
and use the private anchor
\begin{equation}\label{eq:v14-anchor}
 u^2+6v_a^2-2w_a^2-4Cw_a-3T^2=0.
\end{equation}
The inherited quaternion criterion, including parameter zero,
supplies rational witnesses exactly when $T/C$ is integral at2
and3. Its3-integrality forces all three summands in $F_s/C+q$
to be3-integral. Their possible finite valuations lie in the
disjoint classes $2\mathbb Z$, $1+4\mathbb Z$, $3+4\mathbb Z$;
the binary sum uses anisotropy modulo3 and~\eqref{eq:v14-X}.
A negative minimum cannot cancel. Thus $X_A,X_T,X_Z$ are3-units,
so the three root labels are odd. Also $F_s/C$ is3-integral
and is dyadically integral by the scalar lemma. The old unshifted
anchor can therefore be restored with new private witnesses
before H or height. An odd $s$ would make $F_s$ a3-unit,
so $s$ is even and the S-root label $s+1$ is odd.

For an odd root, its half point $[n_i]R_0=(A,B)$ has
$A=(1+\rho)r^2$. Duplication gives
\[
 a_{\rm half}=(A^2+D)/(2B),\qquad
 b_{\rm half}=(B/A)^2/2-c_i,
\]
with $a_{\rm half}^2=c_i$ and the correct product ordinate.
After changing both root signs if necessary, the old half-chart
equations are
$A^2-2(c_i+b_{\rm half})A-D=0$ and
$\operatorname N(A)+\operatorname N(r)^2=0$.
In the final A,T,Z modules the sign freedom is $a_i\mapsto-a_i$,
with reconstructed $b_i\mapsto-b_i$. For S negate
$a_2,b_2,g_0,g_1,\tau$ together. These actions fix the actual
points, $F_s$, all detector values and the private unscaled anchor
witnesses. The scalar rows change at most by an overall sign.

Conversely the old odd roots make $q$ integral at2 and3 by
\eqref{eq:v14-X}, so the same quaternion criterion supplies the
shifted anchor. No prime or H choice changes. The exact private
unscaling used here multiplies the old anchor witnesses and its
parameter $B$ by $g_0$, with inverse multiplication by $\tau$.
It gives $T=g_0B=F_s+Cq$ and~\eqref{eq:v14-anchor}.
This exact step assumes no private anchor witness nonzero.

More explicitly, let the old private residuals be
$R_Q=Q-q$, $R_B=B-2g_1-CQ\tau$, and
$R_A=U^2+6V^2-2W^2-4C\tau W-3B^2$.
Let $N_T=T-2g_0g_1-Cq$ and let $N_A$ be the new
anchor residual. With $G=g_0\tau-1$, pullback by the
forward map gives the unconditional identities
\[
 N_T=g_0R_B+CR_Q+CQG,\qquad
 N_A=g_0^2R_A+4Cg_0WG.
\]
The reverse map is
$Q=q$, $B=\tau T$, $U=\tau u$, $V=\tau v_a$,
$W=\tau w_a$. Its pullbacks are
\[
 R_Q=0,\qquad R_B=\tau N_T+2g_1G,\qquad R_A=\tau^2N_A.
\]
Thus every private anchor consumer is restored with its
literal normalization, and all nonprivate coordinates are fixed.

\subsection{Contact poles and their conditional domains}\label{subsec:v14-contact}
For each scale put
\[
 R_j=F_5(A_j\lambda),\quad E_jR_j=1,\qquad
 F_5(t)=t(t+1/2)(t+1)(t-1/2399)(t+32/407).
\]
The retained contact gates make $R_j$ affine in existing coordinates.
With $W_j=G_{j,0}G_{j,1}G_{j,2}$ retain only
\begin{equation}\label{eq:v14-pole}
 \operatorname N(W_j)=2\operatorname{Re}(W_j)z_j,\qquad
 B_j(R_j+E_j+2)=z_j^2.
\end{equation}
There is no quantified inverse of $\operatorname{Re}(W_j)$.
At a required contact prime, $\alpha=v_p(R_j)>0$ is odd
and $1+R_j$ is a unit. The groups below give $W_j\ne0$.
The norm row then forces its real component and $z_j$ nonzero:
$w_0^2-2w_1^2=0$ over $\mathbb Q$ forces both components zero.
Thus
\[
 T_j^{\rm tr}=\operatorname{Tr}(W_j^{-1})=1/z_j,\qquad
 B_j=\frac{R_j}{(T_j^{\rm tr})^2(1+R_j)^2}.
\]
These domains are needed only at required contact primes.
A detector with radical one has empty local obligations;
no global trace inverse is asserted for it. The case $R_j=-1$
likewise cannot occur at a required contact.

In completeness the ratio between $R_j$ and the earlier contact
rational function is $A_j^5(t-1)^6$, a rational square and a
unit at selected contacts. Hence $R_j$ has valuation one there
and even valuation outside the prescribed odd support.
Both denominators in $B_j$ change valuations by even integers.
This includes additional primes introduced by later witnesses.

\subsection{Common-scale groups and the successful place}\label{subsec:v14-groups}
Use the physical tests
\[
 H_0=f(H),\quad H_i=L_{i+1}-\overline{LP_{i+1}}\ (1\leq i\leq3),
 \quad H_4=\mathcal L(L_2,L_3,L_4),
\]
\[
 M_0=L_1-\overline{LP_1},\quad M_1=L_1-\beta(X),\quad
 P_i=f(U_i)^2,\quad Z_1=\overline{f(V_1)^4},\quad
 T_i=\overline{f(V_i)^8}.
\]
Here $\beta(X)=204(X^2+X)+2555$ is the retained comparison
polynomial. Conjugation is part of the physical tests.
The Lucas polynomial is
\[
 \mathcal L(A,T,Z)=T^2-A(A-2)TZ+(A-2)^2Z^2-1.
\]
The affine recentering $C_L=(A-2)Z$, $V_L=AT$ gives
$\mathcal L+1=T^2+C_L^2-C_LV_L$ without changing its value.

For a nonzero common rational $\gamma$, store
\[
 ZH_i=H_i^4/\gamma^{160},\quad ZM_i=M_i^4/\gamma^{40},\quad
 ZP_2=P_2^2/\gamma^{160},\quad
 ZP_3=P_3^2/\gamma^{80},\quad ZP_4=P_4^2/\gamma^{80},
\]
and $ZP_1=P_1^2$. These definitions are quadratic using the
retained square inputs and common powers. For each detector,
\begin{align}
 G_0&=1+RZH_1+R^2ZH_2+EU+E^4Z_1+\gamma^{16}E^8,\label{eq:v14-g0}\\
 G_1&=1+RZH_3+R^2ZH_4+R^3ZP_2+E^4ZP_4+EZP_1+E^8T_0,\label{eq:v14-g1}\\
 G_2&=1+RZM_0+R^2ZM_1+R^3ZP_3+E^8T_2+E^4T_3+ET_4.\label{eq:v14-g2}
\end{align}
Here $U=H_0^2$ is the actual squared-H output even before
H is reconstructed. The product, trace and pole use these
same final groups.

At a consumed place set $\alpha=v(R)>0$ odd and
$c_\gamma=v(\gamma)\in\mathbb Z$. For a reciprocal point
$h^2=f(1-Df^2)$, a nonzero $f$ has even valuation: the leading
valuation is $v(f)$ when positive and $3v(f)$ when negative.
Zero leaves are omitted. Thus the relevant powers have
periods4,8,16 before formality or logarithms. In the generic
squared-H branch its algebraic reconstruction gives the same
curve fact; its exceptional branches have the fixed signed
values already classified. Other high and middle tests are
literal fourth powers, independently of their later ratio meanings.

In term order, the valuation periods and coefficients of
$\alpha$ are
\[
 \begin{array}{c|l}
 G_0&(4,1),(4,2),(4,-1),(8,-4),(16,-8)\\
 G_1&(4,1),(4,2),(8,3),(8,-4),(8,-1),(16,-8)\\
 G_2&(4,1),(4,2),(8,3),(16,-8),(16,-4),(16,-1).
 \end{array}
\]
All gamma exponents are divisible by the relevant periods.
Within each row the classes modulo16 are disjoint and exclude
zero for every odd $\alpha$. Every group therefore has a unique
minimum valuation, is nonzero, and has valuation at most zero.
This covers arbitrary $c_\gamma$ and every zero-leaf pattern.

Let $\ell_p$ be the maximum valuation of $W$ over places above
$p$. Trace integrality gives $v_p(T^{\rm tr})\geq-\ell_p$.
The pole obstruction and~\eqref{eq:v14-pole} give
\[
 0\leq v_p(B)=\alpha-2v_p(T^{\rm tr})\leq\alpha+2\ell_p.
\]
Its odd parity makes $\ell_p>-\alpha/2$. At a place attaining
the maximum, every group, and then every individual term,
has valuation greater than $-\alpha/2$. This is one place for
all tests, also at an inert unramified prime.

The guard gives $c_\gamma>15\alpha/32$. The high and middle
terms give
\[
 4v(H_i)>160c_\gamma-(r+1/2)\alpha,\qquad
 4v(M_i)>40c_\gamma-(r+1/2)\alpha\quad(r=1\text{ or }2).
\]
For $\alpha=1$ the worst bounds are157.5 and37.5; for
$\alpha=3$ they are312.5 and72.5; for $\alpha\geq5$ they
exceed$72.5\alpha$ and$16.25\alpha$. Thus $v(H_i)\geq40$
and $v(M_i)\geq10$. The primary bounds are
\[
 \begin{split}
 4v(f(U_2))&>160c_\gamma-3.5\alpha,\\
 4v(f(U_3))&>80c_\gamma-3.5\alpha,\\
 4v(f(U_4))&>80c_\gamma+3.5\alpha.
 \end{split}
\]
Using even valuations, these give respectively40,20,22 for the
reciprocal depths. The other low terms force positive valuation
of every required nonzero primary or alternate reciprocal value.

On either exceptional squared-H branch $T_0$ is a unit at every
consumed place. Its $E^8T_0$ term has valuation$-8\alpha$,
contradicting the common strict bound. Either branch would
therefore make all radicals one, contrary to the pre-H
root-height/anchor bounds. The unchanged fixed-point exclusions
from Section~\ref{sec:v13} then restore the full H/V0 domain
before logarithms. Here their required hypothesis is only
positive valuation of $f(H)$, rather than the former stronger
order40 test. A nonidentity torsion value of H, or a value
in either of the fixed translates $\pm P+E[2](K)$, has the
same fixed unit obstruction used there at every consumed
large prime. The new term $EU$, where $U=f(H)^2$ after
reconstruction, gives $2v(f(H))>\alpha/2>0$ and excludes
those fixed values. The identity is already outside the
reconstructed affine H chart. Thus weakening the old H test
does not weaken these domain exclusions; no high-depth
logarithm estimate is used in them.

\subsection{Reciprocal point charts and their orientation}
\label{subsec:v14-reciprocal}
For each of the four root indices use reciprocal coordinates
\[
 f=1/x(U),\quad h=-y(U)/x(U)^2,\qquad
 F=1/x(V),\quad k=-y(V)/x(V)^2.
\]
The pair $(0,0)$ denotes $O$. Retain the square caches $Y=f^2$,
$B=F^2$. In particular the sign of $h$ and $k$ is negative.
The reduced chart uses only $f,h,F,s,Y,B$ and imposes
\begin{align}
 h^2&=f-DfY,\label{eq:v14-primary-curve}\\
 e&=s^2-c+D(f+F)=0,\label{eq:v14-slope-square}\\
 q&=2sh-c(f-F)-1+D(2Y+fF)=0.\label{eq:v14-reduced-slope}
\end{align}
Here $s$ is a private slope, unrelated to the scalar integer
$a-1$ when that integer is denoted by $s$ elsewhere.
All displayed equations have degree at most two in the quantified
coordinates. Products such as $fY$ use the retained cache $Y$.
Define, without adding a coordinate,
\[
 k=s(f-F)-h.
\]
Put $r_1=s(f-F)-h-k$, $r_2=s(h-k)-1+D(Y+fF+B)$ and
$r_4=h^2-f+Df^3$. Exact expansion, using the two caches, gives
\begin{align}
 r_1&=0,\qquad r_2=q-(f-F)e,\label{eq:v14-slope-restoration}\\
 k^2-F+DF^3&=r_4-(f-F)q+(f-F)^2e.\label{eq:v14-alt-restoration}
\end{align}
Thus both reciprocal cubics and the three slope equations hold
before a ratio, valuation or logarithm is interpreted.

These cubics describe actual points of $E(K)$. For $f\ne0$ the
inverse is $x=1/f$, $y=-h/f^2$; for $f=0$ the cubic gives $h=0$
and the point is $O$. This chart omits $(0,0)\in E(K)$ but includes
the other two finite points of order two. Those latter points will
be excluded by the low tests before division by a ratio denominator.
Translation by $(0,0)\in E(K)$ sends the represented point to
$(-Df,-Dh)$, including the case in which the represented point
is $O$.

The oriented difference is
\begin{equation}\label{eq:v14-root-ordinate}
 U-V=(c,y_R),\qquad y_R=Dh-s(Df+c).
\end{equation}
For completeness, this assertion includes every branch of the
equations. If $f\ne F$, the line with slope $s$ joins
$(-Df,-Dh)$ and $(-DF,Dk)$, so the usual secant formula gives
the stated difference. If $f=F$, the first slope equation gives
$k=-h$. The case $h=0$ is impossible: the cubic gives either
$f=0$ or $Df^2=1$, whereas the second slope equation would give
$3Df^2=1$. Hence $h\ne0$, and that equation gives
$2sh=1-3Df^2$, exactly the tangent slope at $(-Df,-Dh)$.
The tangent formula again gives~\eqref{eq:v14-root-ordinate}.
If $U=O$ the equations give $Fc=1$ and $R=-V$; if $V=O$ they
give $fc=1$ and $R=U$. Both points cannot be $O$.

There is also a division-free check on the root curve:
\begin{equation}\label{eq:v14-root-identity}
 y_R^2-c^3+Dc
 =D^2r_4-D(Df+c)(r_2-sr_1)
       +e(c+Df)(c+DF).
\end{equation}
Consequently the retained square equation $a_i^2=c_i$ recovers
the old root chart once $c_i\ne0$: put $b_i=y_R/a_i$, obtaining
$b_i^2=c_i^2-D$.

The nonzero assertion uses no auxiliary ratio. In the literal
system the successive root products are
\[
 c_1c_2=v_{337}+2\rho v_{338},\quad
 c_1c_2c_3=2(v_{339}+2\rho v_{340}),\quad
 c_1c_2c_3c_4=2(v_{341}+2\rho v_{342}).
\]
If any $c_i$ vanished, both last components would vanish.
The unchanged height-product chain would then give
$v_{343}=v_{391}=v_{392}=v_{393}=v_{394}=0$, contradicting its
retained inverse equation. Thus every $c_i$ and $a_i$ is nonzero.
The squareclass calculation now places every root in $2E(K)$.
This supplies the root chart used in the parity and pre-H
arguments above; it does not require the auxiliary points to
lie in $2E(K)$.

For the scalar root retain both old coordinates $a_2,b_2$ and
the old equation $b_2^2=c_2^2-D$. Add the quadratic orientation
equation
\begin{equation}\label{eq:v14-S-orientation}
 a_2b_2=Dh-s(Df+c_2).
\end{equation}
The scalar subsystem consumes this particular ordinate. A root
curve equation alone would allow its negative and would not
identify its scalar point with $U_2-V_2$.
The physical caches for this chart are
$Y=v_{257}+2\rho v_{258}$ and $B=v_{196}+2\rho v_{197}$.
Every old scalar consumer of $a_2,b_2,c_2$ is retained.

Set $N_0=\mathcal S$, the fixed affine expression
in~\eqref{eq:v14-normalizer}. After the H domain is restored
its meaning is $-1/z(H)$. It is not a new quotient or a new
coordinate in the pre-H equations. The primary equations are
\begin{equation}\label{eq:v14-primary-ratios}
 hL+fN_0=0,\qquad (1-DY)L+hN_0=0.
\end{equation}
The retained alternate equation is
\begin{equation}\label{eq:v14-regular-ratio}
 (1-DB)z=s(f-F)-h=k.
\end{equation}
It remains quadratic after substituting for $k$.
At $O$ these equations force $L=0$ or $z=0$, respectively,
rather than allowing a free ratio. The raw alternate equation
$kz=F$ is not imposed. Its recovery requires the global order of
the proof. With $A=1-DB$ and $g=Az-k$,
\begin{equation}\label{eq:v14-raw-ratio}
 A(kz-F)=kg+(k^2-F+DFB).
\end{equation}
The right side vanishes by the curve and regular ratio equations.
One may cancel $A$ only after the low tests have been applied.
At a consumed place they give $v(F)>0$ for a nonidentity point,
so $A$ is a local unit. At the identity $F=0$ and $A=1$.
Thus $A\ne0$ in $K$, and the raw ratio follows.

The order is therefore the same for all four root indices.
First restore the cubics, oriented group differences and nonzero
root charts. Next restore root parity and the scalar-anchor
contradiction to empty contact support. Apply the same-place low
tests using the actual cubics and literal caches. Only then
cancel $1-DB$ and interpret the ratios. In particular the scalar
auxiliary pair is not passed to an older argument which assumed
an unavailable square root of $f$ or $F$.

\subsection{Auxiliary point domains and the H-depth bootstrap}
\label{subsec:v14-logarithms}
We give the domain argument for all four pairs simultaneously.
At a good consumed place $D$ is a unit. On
$h^2=f-Df^3$, a nonzero $f$ has even valuation. If $v(f)>0$,
$2v(h)=v(f)$; if $v(f)<0$, $2v(h)=3v(f)$; and if $v(f)=0$
there is nothing to prove. The identity has $f=0$ and every
positive-depth condition is automatic. Consequently the
valuations of $f^2,f^4,f^8$ have periods $4,8,16$ on all of
$E(K)$ represented by this chart. This proves the periods used
in the group calculation before H has been restored.
The low tests exclude the finite points of order two, whose
reciprocal abscissae are units at these places.

After the fixed H/V0 domain and the raw ratios have been restored,
write
\[
 H=mR_0+T_H,\quad U_i=k_iR_0+T_i,\quad
 V_i=(k_i-2n_i)R_0+T_i,\quad
 V_0=(m-2)R_0+T_H,
\]
where all the $T$'s are in $E[2](K)$. The correct identities,
which kill torsion without imposing a doubled-point domain, are
\begin{equation}\label{eq:v14-doubled-identities}
 [2m]U_i=[2k_i]H,\qquad
 [2(m-2)]V_i=[2(k_i-2n_i)]V_0.
\end{equation}
The restored H exclusions make $H,V_0$ nontorsion.
Taking formal logarithms gives full ratios
\[
 \frac{u_i}{h_0},\qquad \frac{u_i-n_i}{h_0-1},
 \qquad u_i=k_i/2,\quad h_0=m/2.
\]
The half-indices are integral at each consumed odd prime.
At the identity the formal unit
$\log(Q)/z(Q)$ means its regular extension, whose value is one.
No quotient $0/0$ is taken.

One of $h_0,h_0-1$ is a unit. The corresponding full ratio is
integral. The depth-one formal unit estimates and the positive
cross-test transfer integrality to both approximate ratios and
both full ratios. In particular $L_2,L_3,L_4$ are integral at
the common place. Their retained Lucas test has positive
valuation and its value at $(0,0,0)$ is $-1$. At least one of
these three $L_i$ is therefore a unit. The primary bounds proved
above give formal depth at least10 for its point $U_i$.
The exact identity $z(U_i)=L_i z(H)$ then gives depth at least10
for $H$. All integral primary $L_i$ give the same depth for
every nonidentity primary point.

This bootstrap precedes the order40 estimate. At depth10 the
relative primary logarithm error has valuation at least40.
The alternate estimates have order40 for the degree-nine
approximants and order12 for the middle approximants. Their
denominators are units at the consumed place. Since the full
ratios are integral, these relative errors are also absolute
errors. Multiplying the two ratios by $h_0$ and $h_0-1$ and
subtracting gives the retained congruences to the integer root
labels $n_i$, with their required orders10 or40.
No assertion that the global auxiliary labels $u_i$ are integers
is used.

The larger auxiliary domain also requires global nonvanishing
of the rational approximant numerator and denominator.
In its variable $w=Df^2$ their coefficient lists are
\[
\begin{split}
N(w):\;&1,-9776/5365,2941/3219,-1024/13949,\\
 &-1920/237133,-1024/711399,-1408/5928325,0,
 2432/103152855,\\
Q(w):\;&1,-1526/1073,513/1073 .
\end{split}
\]
The unchanged certificates modulo71 for $N$ and modulo7 for $Q$
apply to every element of $K$. Both primes split in $K$; the
respective coefficients are integral, with unit leading
coefficient, and their reductions have no root. A root with
negative valuation is excluded by its unique leading term,
and an integral root is excluded by reduction. Thus the old
nonzero products and inverses remain justified. A positivity
interval restricted to doubled real points is not used here.

The retained root-height, Lucas and scalar-depth arguments can
now be applied with their original inputs. The roots are the
same odd integer multiples, the congruences have the required
precision, the scalar point has the orientation
in~\eqref{eq:v14-S-orientation}, and the scalar depth24 transfer
is unchanged. They conclude that the distinguished input is
an integer. This proves soundness of the changed arithmetic
certificates.

\subsection{Simultaneous completeness of the common scale}
\label{subsec:v14-completeness}
We check that a single choice of witnesses meets all the new
tests. Start with the original simultaneous construction for
an integer input, with the same positive progression and the
same contact-prime restrictions. Use its depth19 preparation.
The resulting primary formal depths are at least20, the high
tests have valuation at least40, and the middle tests have
valuation at least12. The fixed H/V0 exclusions remain part of
this construction.

At each contact prime $\alpha=v(R)=1$. Prescribe $v(\gamma)=1$.
In the three groups the high-term valuations are at least1
or2 and the middle-term valuations at least9 or10. The primary
terms have lower bounds159,3,83,76 in root order; the guard has
lower bound8. The alternate low terms are unchanged and have
positive valuation. Every group is therefore one plus an
element of the maximal ideal at the successful place.
At the other place above a split contact prime, the inverse
of each group has residue zero or one. Thus the trace of the
product of inverses has residue one or two and is a unit.
The pole value has positive odd valuation one.

At each preparation prime $p=2,3,17$, put
$\eta_j=v_p(A_j)$, prescribe $v_p(\gamma)=v_p(\lambda)=-\nu_p$,
and enlarge $\nu_p$ before choosing the contact primes.
Since $v_p(R)=5(\eta_j-\nu_p)$, the coefficients multiplying
the bounded test values have valuations
\[
\begin{array}{c|c}
 \gamma^{-160}R&155\nu_p+5\eta_j\\
 \gamma^{-160}R^2&150\nu_p+10\eta_j\\
 \gamma^{-160}R^3&145\nu_p+15\eta_j\\
 \gamma^{-40}R&35\nu_p+5\eta_j\\
 \gamma^{-40}R^2&30\nu_p+10\eta_j\\
 \gamma^{-80}R^3&65\nu_p+15\eta_j\\
 \gamma^{-80}E^4&100\nu_p-20\eta_j\\
 \gamma^{16}E^8&24\nu_p-40\eta_j .
\end{array}
\]
Each tends to $+\infty$ with $\nu_p$. These are all the new
coefficient forms. The previous finite preparation therefore
works after increasing its bounds.
At the real places impose $\lambda<\gamma<2\lambda$.
The corresponding powers of $\lambda$ are
$-155,-150,-145,-35,-30,-65,-100,-24$.
They all tend to zero uniformly in this interval. The groups
tend to one, their trace to two, and the pole value tends to
$1/(4R)$ through positive values. This gives the real conditions
required by the unchanged local-global pole certificates.

Away from the contact support, $R$ has even valuation. Since
the other factors in its pole value are squares, new primes in
the denominator of $\gamma$ create no odd pole support.
This also explains why $\gamma$ may be chosen on the anchor
quadric below rather than independently. The argument requires
its specified local valuations and real interval, not an
absence of other denominator primes.

The choices have a fixed order. First prepare the scalar and
root witnesses and all finite and real bounds. Choose the
contact primes and then $\lambda$. Choose the anchor point
and hence $\gamma$ with the required local conditions.
Finally choose H in the original progression, deep enough at
the finitely many selected places and in the retained real
neighborhoods. The bounds above are uniform in the anchor
choice. Thus this last H choice meets all three groups
simultaneously.

\subsection{Coupling the scale to the quaternion anchor}
\label{subsec:v14-anchor}
Let $T$ be the shifted scalar anchor already constructed.
Writing the old quaternion coordinates as $u,v,w$, put
$z=w+C$. Its equation is
\begin{equation}\label{eq:v14-anchor-quadric}
 u^2+6v^2-2z^2=N,\qquad N=3T^2-2C^2.
\end{equation}
Here $N\ne0$, since $N=0$ would make the rational square
$(T/C)^2$ equal $2/3$.
Use $u$ to define $\gamma=1/u$. The independent common-scale
chain is replaced by the eight coordinates specified by
\[
 a_2=u^2,\quad a_4=a_2^2,\quad a_8=a_4^2,\quad
 a_{16}=a_8^2,\quad a_{16}J=1,\quad
 g_{40}a_8=J^2,\quad g_{80}=g_{40}^2,\quad
 g_{160}=g_{80}^2.
\]
These equations force $u\ne0$, $J=\gamma^{16}$ and
$g_j=\gamma^j$ for $j=40,80,160$. They recover the former
common-scale coordinates uniquely and prove soundness of
this replacement.

We prove that the extra choice of $u$ does not obstruct
completeness. In the original scalar preparation require also
the allowed condition $F_s\in8\mathbb Z_2$. The same positive
progression and further formal depths remain available.
For the three odd roots $v_2(X_i)=-1$, and
\[
 q=3(2X_A)^2+3(2X_T)^2+27(2X_Z)^2\equiv1\pmod8.
\]
The shifted anchor therefore has $t=T/C$ odd at2 and integral
at3. The retained quaternion criterion supplies one rational
point $P$ on~\eqref{eq:v14-anchor-quadric}.
Fix it before choosing the contact primes.

The required local sets on this quadric are nonempty. Choose
$\nu_2,\nu_{17}\geq1$ and
$\nu_3>v_3(C)=2^{25}$, enlarged to meet all earlier preparation
bounds. At2 take
\[
 u=2^{\nu_2},\qquad v=T/2,\qquad
 z=(C/2)\sqrt{4-3t^2+2(u/C)^2}.
\]
The radicand is one modulo8. At3 take
\[
 u=3^{\nu_3},\quad v=0,\quad
 z=C\sqrt{1-3t^2/2+(u/C)^2/2};
\]
the radicand is one modulo3. At17 take
\[
 u=17^{\nu_{17}},\quad v=T/\sqrt2,\quad
 z=C\sqrt{1+u^2/(2C^2)}.
\]
Here2 is a square in $\mathbb Q_{17}$ and the last radicand
is one modulo17. Each construction solves the quadric exactly.

At a contact prime $p>76176$ exclude the finitely many
denominator primes of $T$, and prescribe $u=p^{-1}$.
Setting $V=pv$, $Z=pz$ leaves
$6V^2-2Z^2=p^2N-1$.
Modulo $p$ the right side is $-1$. The binary form
$Z^2-3V^2$ represents every nonzero residue: it is a split
product if3 is a square and a surjective finite-field norm
otherwise. The resulting solution is nonsingular, so Hensel
lifting gives the desired local point. At the real place take
$1/(2\lambda)<u<1/\lambda$, choose $v$ sufficiently large,
and solve for a positive real $z$. Hence $\lambda<\gamma<2\lambda$.

For clarity, weak approximation on this particular quadric
can be proved without adding quantified witnesses.
Let $Q(u,v,z)=u^2+6v^2-2z^2$ and let $B$ be its associated
symmetric bilinear form, so $B(X,X)=Q(X)$.
For a direction $d$ with $Q(d)\ne0$ put
\[
 \Phi(d)=P-\frac{2B(P,d)}{Q(d)}d .
\]
Then $Q(\Phi(d))=N$. Every point $X$ on the quadric with
$Q(X-P)\ne0$ is obtained by taking $d=X-P$.
This image meets every nonempty local open set. Indeed the
excluded tangent section has equation $B(P,X)=N$.
If $P$ and $X$ are not proportional, choose
$h\perp X$ with $B(P,h)\ne0$ and use the local curve
\[
 X(s)=\sqrt{1-s^2Q(h)/N}\,X+sh .
\]
It stays on the quadric and leaves that section to first
order. If $P$ is proportional to $X$ on the section, then
$P=X$. The tangent plane is nondegenerate because $N\ne0$;
choose $h$ there with $Q(h)\ne0$. The same curve leaves the
section to second order. The square root exists for
sufficiently small $s$, also over $\mathbb Q_2$.

Finally approximate the three direction coordinates
simultaneously by rationals. After clearing the local
denominators with an integer $D_0$, impose the required finite
congruences modulo an integer $M$. For integers
$L\equiv1\pmod M$ tending to infinity, coordinates of the form
$(A+Mk)/(D_0L)$ retain those congruences and have real mesh
$M/(D_0L)$ tending to zero. The Chinese remainder theorem
chooses $A$ for each coordinate. Continuity of $\Phi$ on the
chosen local neighborhoods now supplies a rational point in
all the required open sets. Its $u$ gives precisely the
valuations and real interval required for $\gamma$.
This proves simultaneous completeness of the coupled anchor.

\subsection{Exact private changes and the finite polynomial}
\label{subsec:v14-finite}
We relate the preceding arithmetic argument to the literal
quadratic presentation. Coordinates represent elements of
$K$ by their rational components; each displayed field
equation contributes its two component equations.
All rational coefficients are cleared and each equation is
divided by its integer content before taking the sum of
squares. These operations do not change its zero set.

For the reciprocal replacement, an old root chart with
scaled private coordinates $V_{\rm old}$, $I_{\rm old}$,
$Q_{\rm old}$, $K_{\rm old}$ has $c=a_i^2\ne0$.
Its actual reciprocal coordinates are
\[
 f=V_{\rm old}^2/c,\quad
 h=-I_{\rm old}V_{\rm old}/c,\quad
 F=Q_{\rm old}^2/c,\quad
 k=-K_{\rm old}Q_{\rm old}/c.
\]
The old identities $I_{\rm old}^2=c(1-DY)$ and
$K_{\rm old}^2=c(1-DB)$, together with
$I_{\rm old}L=V_{\rm old}N_0$ and
$K_{\rm old}z=-Q_{\rm old}$, give the new curve and both
regular ratio equations, including the identity cases.
The actual equality $U-V=Q_i$ supplies a rational secant
slope or its rational tangent limit. For the scalar pair it
also supplies~\eqref{eq:v14-S-orientation}.
No new local neighborhood, root sign, anchor choice or H
choice is needed for this change.
Conversely~\eqref{eq:v14-root-identity} restores the deleted
root ordinates and the preceding soundness proof applies to
the enlarged auxiliary domain. We assert equality of input
projections, not a bijection to all the old private
square-root fibers.

The subsequent deletion of each alternate ordinate is an
exact change on the full zero set. Its inverse is
$k=s(f-F)-h$, and~\eqref{eq:v14-alt-restoration} and the
globally justified~\eqref{eq:v14-raw-ratio} recover every
deleted equation. This accounts for six rational coordinates
in the A, T and Z pairs. The scalar pair saves two further
coordinates while retaining its root ordinate and adding
the orientation equation.

The remaining private replacements are exact power chains
or polynomial changes. For a local inverse $ER=1$, a shorter
chain may specify
\[
\begin{gathered}
 EJ=R^2,\quad E_4J=E,\quad E_8=E_4^2,\quad E_{19}J=E_8^2,\\
 RE_{39}=E_{19}^2,\quad RE_{79}=E_{39}^2,\quad
 E_{158}=E_{79}^2,\quad E_{157}=RE_{158}.
\end{gathered}
\]
It forces $J=R^3$ and exactly the indicated powers of $E$.
There is no zero or sign branch because $ER=1$.
The final common-scale presentation needs only the local
coordinates for $E,R^2,R^3,E^4,E^8$, with shared powers
specified above. Its equivalence as an arithmetic
certificate was proved in the common-place and completeness
arguments, rather than inferred from coordinate counting.

The sparse Lucas change replaces its private coordinate
$C_{\rm old}=AZ$ by $C_{\rm new}=C_{\rm old}-2Z$.
The inverse is $C_{\rm old}=C_{\rm new}+2Z$, and the defining
polynomial becomes
\[
 F+1=T^2+C_{\rm new}^2-C_{\rm new}(AT).
\]
Both substitutions are affine in the private coordinates.
The cached product $AT$ makes the displayed equation
quadratic. Every external Lucas value is unchanged.
This change reduces terms, not the number of witnesses.

For reference the successive total-variable counts of the
selected construction are
\[
394,\ 392,\ 390,\ 372,\ 369,\ 366,\ 326,\ 319,\
318,\ 315,\ 314,\ 313,\ 307,\ 301,\ 299 .
\]
The early exact eliminations, affine H chart and squared-H
replacement are justified above. The count326 replaces
separate pole blocks by the three groups; the final group
argument proves their required arithmetic properties
directly. The count315 precedes common scaling.
Common scaling removes twelve local coordinates, adds nine
shared coordinates and two scaled test coordinates, giving314.
The coupled quaternion anchor replaces those nine by eight,
giving313. Reciprocal A, T and Z root charts give307,
alternate-ordinate elimination gives301, and the oriented
scalar auxiliary chart gives299.
Intermediate alternatives with different depth requirements
are not needed for this selected proof.

The literal final system has279 quadratic component
equations in299 active rational variables, including the
input. Its polynomial is their unweighted sum of squares.
Over $\mathbb Q$ it vanishes exactly when every component
equation vanishes. Exact expansion, collection by exponent
vector and removal of zero coefficients give6,630 monomials,
total degree four, content one, and maximum absolute
coefficient $186624C^2=186624\cdot3^{2^{26}}$.
Every coordinate is active. The coefficient artifacts specify
the exponent-vector order and the integer coefficient
representation, so these finite assertions refer to a
single polynomial, not just to its zero set.

\subsection{The ordinary coefficient presentation}
\label{subsec:v14-ordinary}
The large fixed coefficient can be represented with fourteen
constant coordinates. Put
\[
 (a_0,\ldots,a_{13})=(2049,
 \underbrace{2048,\ldots,2048}_{11\ {\rm entries}},0,0)
\]
and impose
$c_0=3^{a_0}$ and $c_j=3^{a_j}c_{j-1}^2$ for $1\leq j\leq13$.
If $c_j=3^{e_j}$ then
$e_0=2049$ and $e_j=2e_{j-1}+a_j$.
After the eleven2048 steps $e_{11}=2^{23}$, and the two
last squares give $e_{13}=2^{25}$. Thus $c_{13}=C$.
Each coordinate is uniquely determined; no square-root sign
is introduced.

Add one more coordinate
\[
 Q=48v_{83}^2+48v_{87}^2+432v_{89}^2.
\]
Only two rows of the symbolic system contain $C$. Replace
them by
\[
\begin{split}
 v_{55}^2+6v_{56}^2-2v_{57}^2
       -4c_{13}v_{57}-3v_{903}^2&=0,\\
 v_{903}-2v_{543}v_{544}-c_{13}Q&=0 .
\end{split}
\]
Both are quadratic. Substitution of the unique constants
and $Q$ gives the exact symbolic rows. Projection and unique
extension therefore give inverse maps on full rational
solution sets and fix the input.
There are $299+14+1=314$ variables and $279+14+1=294$
quadratic equations.

Apply the same primitive integer normalization and
unweighted sum of squares. The ordinary polynomial has
6,672 nonzero monomials and content one.
Let $K_0$ denote the constant coefficient of the symbolic
sum of squares, an explicitly tabulated positive integer
of1,514 bits. The exact largest absolute coefficient of
the ordinary expansion is
\[
 3^{4098}+K_0,\qquad
 2^{6495}\leq3^{4098}+K_0<2^{6496}.
\]
This proves the asserted ordinary coefficient bound.
The source-coordinate table retains the original coordinate
identifiers and has222,890 bytes. A compact table renames
the active variables consecutively by an explicit bijection
fixing the input; its byte count is a separate artifact
property. That renaming changes neither the equations nor
the coefficient statistics. The proof and the coefficient
tables above establish Theorem~\ref{thm:v14-quartic}.
